\section{Mathematical Proofs}
\label{app:proofs}

This section provides detailed mathematical proofs and derivations for key components of our method.

\subsection{KL Divergence with Mixture Prior}

\subsubsection{Definition and Computation}

\textbf{Theorem 1} (KL Divergence with Mixture Prior): The KL divergence between a posterior distribution $q_i(z) = \mathcal{N}(z; \mu_i, \sigma_i^2 I)$ and a mixture prior $p_{\text{mixture}}(z) = \sum_{j=1}^{|\mathcal{S}|} w_j \cdot \mathcal{N}(z; \mu_j, \sigma_j^2 I)$ is given by:

\begin{equation}
\mathbb{D}_{KL}(q_i \,\|\, p_{\text{mixture}}) = \mathbb{E}_{z \sim q_i} \left[ \log q_i(z) - \log p_{\text{mixture}}(z) \right],
\end{equation}

where $\mathcal{S}$ is the set of participating clients from the previous round and $w_j$ are normalized weights (typically sample-size weighted: $w_j = n_j / \sum_{k \in \mathcal{S}} n_k$).

\textbf{Proof:} By definition of KL divergence:

\begin{align}
\mathbb{D}_{KL}(q_i \,\|\, p_{\text{mixture}}) &= \int q_i(z) \log \frac{q_i(z)}{p_{\text{mixture}}(z)} dz \nonumber \\
&= \int q_i(z) \left[ \log q_i(z) - \log p_{\text{mixture}}(z) \right] dz \nonumber \\
&= \mathbb{E}_{z \sim q_i} \left[ \log q_i(z) - \log p_{\text{mixture}}(z) \right].
\end{align}

For a multivariate Gaussian posterior $q_i(z) = \mathcal{N}(z; \mu_i, \sigma_i^2 I)$ with dimension $d$, we have:

\begin{equation}
\log q_i(z) = -\frac{d}{2}\log(2\pi) - \frac{1}{2}\sum_{j=1}^{d} \log\sigma_{i,j}^2 - \frac{1}{2}\sum_{j=1}^{d} \frac{(z_j - \mu_{i,j})^2}{\sigma_{i,j}^2}.
\end{equation}

For the mixture prior, we compute:

\begin{equation}
\log p_{\text{mixture}}(z) = \log \left[ \sum_{j=1}^{|\mathcal{S}|} w_j \cdot \mathcal{N}(z; \mu_j, \sigma_j^2 I) \right].
\end{equation}

Using Monte Carlo estimation with batch size $B$:

\begin{equation}
\mathbb{D}_{KL} \approx \frac{1}{B} \sum_{b=1}^{B} \left[ \log q_i(z^{(b)}) - \log p_{\text{mixture}}(z^{(b)}) \right],
\end{equation}

where $z^{(b)} \sim q_i$ is sampled using the reparameterization trick:

\begin{equation}
z^{(b)} = \mu_i + \sigma_i \odot \epsilon^{(b)}, \quad \epsilon^{(b)} \sim \mathcal{N}(0, I),
\end{equation}

and $\sigma_i = \exp(0.5 \cdot \log\sigma_i^2)$ with $\odot$ denoting element-wise multiplication.

$\square$

\subsubsection{Log-Sum-Exp Trick for Numerical Stability}

\textbf{Theorem 2} (Log-Sum-Exp Trick): For numerical stability when computing $\log p_{\text{mixture}}(z)$, we use the log-sum-exp trick:

\begin{equation}
\log\left(\sum_{i=1}^{N} \exp(a_i)\right) = \max_i a_i + \log\left(\sum_{i=1}^{N} \exp(a_i - \max_i a_i)\right).
\end{equation}

\textbf{Proof:} We factor out the maximum term:

\begin{align}
\sum_{i=1}^{N} \exp(a_i) &= \exp(\max_i a_i) \cdot \sum_{i=1}^{N} \exp(a_i - \max_i a_i).
\end{align}

Taking the logarithm of both sides:

\begin{equation}
\log\left(\sum_{i=1}^{N} \exp(a_i)\right) = \max_i a_i + \log\left(\sum_{i=1}^{N} \exp(a_i - \max_i a_i)\right).
\end{equation}

Since $\exp(a_i - \max_i a_i) \in [0, 1]$ for all $i$, this formulation is numerically stable and prevents overflow/underflow.

For the mixture prior, we apply this trick with:

\begin{equation}
a_j = \log w_j + \log \mathcal{N}(z; \mu_j, \sigma_j^2 I),
\end{equation}

where:

\begin{equation}
\log \mathcal{N}(z; \mu_j, \sigma_j^2 I) = -\frac{d}{2}\log(2\pi) - \frac{1}{2}\sum_{k=1}^{d} \log\sigma_{j,k}^2 - \frac{1}{2}\sum_{k=1}^{d} \frac{(z_k - \mu_{j,k})^2}{\sigma_{j,k}^2}.
\end{equation}

$\square$

\subsection{Reparameterization Trick and Gradient Flow}

\textbf{Theorem 3} (Reparameterization Trick Gradient Flow): Using the reparameterization trick, gradients with respect to $z$ flow to $\mu$ and $\log\sigma^2$.

\textbf{Proof:} The reparameterization trick expresses the random variable $z$ as a deterministic function of parameters and noise:

\begin{equation}
z = \mu + \epsilon \odot \exp(0.5 \cdot \log\sigma^2), \quad \epsilon \sim \mathcal{N}(0, I),
\end{equation}

where $\sigma = \exp(0.5 \cdot \log\sigma^2)$.

The partial derivatives are:

\begin{align}
\frac{\partial z}{\partial \mu} &= I \quad \text{(identity matrix)}, \\
\frac{\partial z}{\partial \log\sigma^2} &= \epsilon \odot \exp(0.5 \cdot \log\sigma^2) \odot 0.5.
\end{align}

By the chain rule, for any function $f(z)$:

\begin{align}
\frac{\partial f(z)}{\partial \mu} &= \frac{\partial f(z)}{\partial z} \cdot \frac{\partial z}{\partial \mu} = \frac{\partial f(z)}{\partial z}, \\
\frac{\partial f(z)}{\partial \log\sigma^2} &= \frac{\partial f(z)}{\partial z} \cdot \frac{\partial z}{\partial \log\sigma^2} = \frac{\partial f(z)}{\partial z} \cdot \epsilon \odot \sigma \odot 0.5.
\end{align}

This ensures that gradients can flow through the sampling operation, enabling end-to-end training of the variational encoder.

$\square$

\subsection{Orthogonal Loss Formulation}

\subsubsection{Pull Loss}

The pull loss encourages latent representations $z$ to align with their assigned prototypes:

\begin{equation}
\mathcal{L}_{\text{pull}}(z) = \|z - \mathbf{p}_{y_i^*}\|_2^2,
\end{equation}

where $\mathbf{p}_{y_i^*}$ is the prototype assigned to client $i$ based on the server's orthogonal label $y_i^* \in \{0, 1\}$.

\subsubsection{Orthonormality Constraint}

To maintain orthogonality between prototypes, we enforce an orthonormality constraint:

\begin{equation}
\mathcal{L}_{\text{orthonorm}} = \|\mathbf{P}^T \mathbf{P} - \mathbf{I}\|_F^2,
\end{equation}

where $\mathbf{P} = [\mathbf{p}_0, \mathbf{p}_1]^T$ is the prototype matrix and $\|\cdot\|_F$ denotes the Frobenius norm.

This constraint ensures that $\mathbf{p}_0^T \mathbf{p}_1 = 0$ (orthogonality) and $\|\mathbf{p}_0\|_2 = \|\mathbf{p}_1\|_2 = 1$ (normalization).

\subsubsection{Total Orthogonal Loss}

The combined orthogonal loss is:

\begin{equation}
\mathcal{L}_{\text{orthogonal}}(z) = \lambda \cdot \mathcal{L}_{\text{pull}}(z) + \gamma \cdot \mathcal{L}_{\text{orthonorm}},
\end{equation}

where $\lambda = 1.0$ (orthogonal weight) and $\gamma = 0.1$ (orthonorm weight) are hyperparameters.

This loss encourages latent representations to cluster around their assigned prototypes while ensuring that different preference types occupy orthogonal subspaces, thereby suppressing general features that are not preference-specific and preventing neural collapse.

\subsection{Evidence Lower Bound (ELBO) Derivation}

\textbf{Theorem 4} (ELBO Derivation): The Evidence Lower Bound (ELBO) for variational inference is:

\begin{equation}
\mathcal{L}_{\text{ELBO}} = \mathbb{E}_{q_\phi(z \mid \mathcal{D}_i)} \left[ \log p_\theta(\mathcal{D}_i \mid z) \right] - \beta \,\mathbb{D}_{KL}(q_\phi(z \mid \mathcal{D}_i) \,\|\, p(z)),
\end{equation}

where $\beta$ is a regularization coefficient.

\textbf{Proof:} We start with the log-likelihood of the data:

\begin{equation}
\log p_\theta(\mathcal{D}_i) = \log \int p_\theta(\mathcal{D}_i \mid z) p(z) dz.
\end{equation}

Introducing the variational posterior $q_\phi(z \mid \mathcal{D}_i)$:

\begin{align}
\log p_\theta(\mathcal{D}_i) &= \log \int q_\phi(z \mid \mathcal{D}_i) \cdot \frac{p_\theta(\mathcal{D}_i \mid z) p(z)}{q_\phi(z \mid \mathcal{D}_i)} dz \nonumber \\
&= \log \mathbb{E}_{q_\phi(z \mid \mathcal{D}_i)} \left[ \frac{p_\theta(\mathcal{D}_i \mid z) p(z)}{q_\phi(z \mid \mathcal{D}_i)} \right].
\end{align}

Applying Jensen's inequality (since $\log$ is concave):

\begin{align}
\log p_\theta(\mathcal{D}_i) &\geq \mathbb{E}_{q_\phi(z \mid \mathcal{D}_i)} \left[ \log \frac{p_\theta(\mathcal{D}_i \mid z) p(z)}{q_\phi(z \mid \mathcal{D}_i)} \right] \nonumber \\
&= \mathbb{E}_{q_\phi(z \mid \mathcal{D}_i)} \left[ \log p_\theta(\mathcal{D}_i \mid z) \right] + \mathbb{E}_{q_\phi(z \mid \mathcal{D}_i)} \left[ \log \frac{p(z)}{q_\phi(z \mid \mathcal{D}_i)} \right] \nonumber \\
&= \mathbb{E}_{q_\phi(z \mid \mathcal{D}_i)} \left[ \log p_\theta(\mathcal{D}_i \mid z) \right] - \mathbb{D}_{KL}(q_\phi(z \mid \mathcal{D}_i) \,\|\, p(z)).
\end{align}

Introducing the regularization coefficient $\beta$ to control the strength of the KL regularization:

\begin{equation}
\mathcal{L}_{\text{ELBO}} = \mathbb{E}_{q_\phi(z \mid \mathcal{D}_i)} \left[ \log p_\theta(\mathcal{D}_i \mid z) \right] - \beta \,\mathbb{D}_{KL}(q_\phi(z \mid \mathcal{D}_i) \,\|\, p(z)).
\end{equation}

The first term is the reconstruction loss (preference alignment), and the second term is the regularization (prior matching). Maximizing the ELBO is equivalent to minimizing the negative ELBO:

\begin{equation}
\mathcal{L}_i(\theta, \phi) = - \text{ELBO}(\mathcal{D}_i) = - \mathbb{E}_{q_\phi} \left[ \log p_\theta(\mathcal{D}_i \mid z) \right] + \beta \,\mathbb{D}_{KL}(q_\phi(z \mid \mathcal{D}_i) \,\|\, p(z)).
\end{equation}

$\square$

\subsection{Standard Gaussian Prior KL Divergence}

For the baseline VPL ablation, we use a standard Gaussian prior $p(z) = \mathcal{N}(0, I)$. The KL divergence has a closed-form expression:

\textbf{Theorem 5} (Standard Gaussian Prior KL): For $q_i(z) = \mathcal{N}(z; \mu_i, \sigma_i^2 I)$ and $p(z) = \mathcal{N}(0, I)$, the KL divergence is:

\begin{equation}
\mathbb{D}_{KL}(q_i \,\|\, \mathcal{N}_0) = \frac{1}{2} \sum_{j=1}^d \left( \mu_{i,j}^2 + \sigma_{i,j}^2 - \log \sigma_{i,j}^2 - 1 \right).
\end{equation}

\textbf{Proof:} For two multivariate Gaussians, the KL divergence is:

\begin{align}
\mathbb{D}_{KL}(\mathcal{N}(\mu_1, \Sigma_1) \,\|\, \mathcal{N}(\mu_2, \Sigma_2)) &= \frac{1}{2} \left[ \text{tr}(\Sigma_2^{-1} \Sigma_1) + (\mu_2 - \mu_1)^T \Sigma_2^{-1} (\mu_2 - \mu_1) - d + \log \frac{|\Sigma_2|}{|\Sigma_1|} \right].
\end{align}

For $q_i = \mathcal{N}(\mu_i, \sigma_i^2 I)$ and $p = \mathcal{N}(0, I)$:

\begin{align}
\mathbb{D}_{KL}(q_i \,\|\, \mathcal{N}_0) &= \frac{1}{2} \left[ \text{tr}(I^{-1} \cdot \sigma_i^2 I) + (0 - \mu_i)^T I^{-1} (0 - \mu_i) - d + \log \frac{|I|}{|\sigma_i^2 I|} \right] \nonumber \\
&= \frac{1}{2} \left[ \text{tr}(\sigma_i^2 I) + \mu_i^T \mu_i - d - \log |\sigma_i^2 I| \right] \nonumber \\
&= \frac{1}{2} \left[ \sum_{j=1}^d \sigma_{i,j}^2 + \sum_{j=1}^d \mu_{i,j}^2 - d - \sum_{j=1}^d \log \sigma_{i,j}^2 \right] \nonumber \\
&= \frac{1}{2} \sum_{j=1}^d \left( \mu_{i,j}^2 + \sigma_{i,j}^2 - \log \sigma_{i,j}^2 - 1 \right).
\end{align}

$\square$

\subsection{Gumbel-Softmax for Differentiable Prior Sampling}

For sampling from the mixture prior (used in visualization and generation), we employ Gumbel-Softmax relaxation \citep{jang2016categorical} with temperature $\tau = 1.0$ to enable differentiable sampling.

\textbf{Component probabilities:}

\begin{equation}
\alpha_j = \frac{\exp((\log w_j + g_j) / \tau)}{\sum_{k=1}^{|\mathcal{S}|} \exp((\log w_k + g_k) / \tau)},
\end{equation}

where $g_j \sim \text{Gumbel}(0, 1)$ are independent Gumbel random variables.

\textbf{Sampling:}

\begin{equation}
z_{\text{prior}} = \sum_{j=1}^{|\mathcal{S}|} \alpha_j \cdot z_j, \quad \text{where } z_j \sim \mathcal{N}(\mu_j, \sigma_j^2 I).
\end{equation}

As $\tau \to 0$, this approaches categorical sampling (hard assignment), while $\tau > 0$ provides a smooth, differentiable approximation.

\subsection{Gradient Computation for Mixture Prior}

\textbf{Theorem 6} (Gradient of KL with Respect to Posterior Parameters): The gradient of the KL divergence with respect to posterior parameters $\mu_i$ and $\log\sigma_i^2$ is:

\begin{align}
\frac{\partial \mathbb{D}_{KL}}{\partial \mu_i} &= \frac{1}{B} \cdot \left( \frac{z - \mu_i}{\sigma_i^2} + \sum_{j=1}^{|\mathcal{S}|} \text{softmax}_j \cdot \frac{z - \mu_j}{\sigma_j^2} \right), \\
\frac{\partial \mathbb{D}_{KL}}{\partial \log\sigma_i^2} &= \frac{1}{B} \cdot \left( -\frac{1}{2} + \frac{(z - \mu_i)^2}{2\sigma_i^2} + \sum_{j=1}^{|\mathcal{S}|} \text{softmax}_j \cdot \frac{z - \mu_j}{\sigma_j^2} \cdot \epsilon \odot \sigma_i \odot 0.5 \right),
\end{align}

where $\text{softmax}_j = \exp(a_j - \max_k a_k) / \sum_{k=1}^{|\mathcal{S}|} \exp(a_k - \max_k a_k)$ with $a_j = \log w_j + \log \mathcal{N}(z; \mu_j, \sigma_j^2 I)$.

\textbf{Proof:} The KL divergence is $\mathbb{D}_{KL} = \frac{1}{B} \sum_{b=1}^{B} [\log q_i(z^{(b)}) - \log p_{\text{mixture}}(z^{(b)})]$, where $z^{(b)} = \mu_i + \sigma_i \odot \epsilon^{(b)}$.

The gradient has two paths:

\textbf{Path 1:} Through $\log q_i(z)$:

\begin{align}
\frac{\partial \mathbb{D}_{KL}}{\partial \log q_i(z)} &= \frac{1}{B}, \\
\frac{\partial \log q_i(z)}{\partial \mu_i} &= \frac{z - \mu_i}{\sigma_i^2}, \\
\frac{\partial \log q_i(z)}{\partial \log\sigma_i^2} &= -\frac{1}{2} + \frac{(z - \mu_i)^2}{2\sigma_i^2}.
\end{align}

By chain rule:
\begin{align}
\frac{\partial \mathbb{D}_{KL}}{\partial \mu_i} \Big|_{\text{Path 1}} &= \frac{1}{B} \cdot \frac{z - \mu_i}{\sigma_i^2}, \\
\frac{\partial \mathbb{D}_{KL}}{\partial \log\sigma_i^2} \Big|_{\text{Path 1}} &= \frac{1}{B} \cdot \left( -\frac{1}{2} + \frac{(z - \mu_i)^2}{2\sigma_i^2} \right).
\end{align}

\textbf{Path 2:} Through $\log p_{\text{mixture}}(z)$:

\begin{align}
\frac{\partial \mathbb{D}_{KL}}{\partial \log p_{\text{mixture}}(z)} &= -\frac{1}{B}, \\
\frac{\partial \log p_{\text{mixture}}(z)}{\partial z} &= \sum_{j=1}^{|\mathcal{S}|} \text{softmax}_j \cdot \frac{-(z - \mu_j)}{\sigma_j^2} = -\sum_{j=1}^{|\mathcal{S}|} \text{softmax}_j \cdot \frac{z - \mu_j}{\sigma_j^2}.
\end{align}

By the reparameterization trick:
\begin{align}
\frac{\partial z}{\partial \mu_i} &= I, \\
\frac{\partial z}{\partial \log\sigma_i^2} &= \epsilon \odot \sigma_i \odot 0.5.
\end{align}

By chain rule:
\begin{align}
\frac{\partial \mathbb{D}_{KL}}{\partial z} &= -\frac{1}{B} \cdot \left( -\sum_{j=1}^{|\mathcal{S}|} \text{softmax}_j \cdot \frac{z - \mu_j}{\sigma_j^2} \right) = \frac{1}{B} \cdot \sum_{j=1}^{|\mathcal{S}|} \text{softmax}_j \cdot \frac{z - \mu_j}{\sigma_j^2}, \\
\frac{\partial \mathbb{D}_{KL}}{\partial \mu_i} \Big|_{\text{Path 2}} &= \frac{\partial \mathbb{D}_{KL}}{\partial z} \cdot \frac{\partial z}{\partial \mu_i} = \frac{1}{B} \cdot \sum_{j=1}^{|\mathcal{S}|} \text{softmax}_j \cdot \frac{z - \mu_j}{\sigma_j^2}, \\
\frac{\partial \mathbb{D}_{KL}}{\partial \log\sigma_i^2} \Big|_{\text{Path 2}} &= \frac{\partial \mathbb{D}_{KL}}{\partial z} \cdot \frac{\partial z}{\partial \log\sigma_i^2} = \frac{1}{B} \cdot \sum_{j=1}^{|\mathcal{S}|} \text{softmax}_j \cdot \frac{z - \mu_j}{\sigma_j^2} \cdot \epsilon \odot \sigma_i \odot 0.5.
\end{align}

Combining both paths:
\begin{align}
\frac{\partial \mathbb{D}_{KL}}{\partial \mu_i} &= \frac{1}{B} \cdot \frac{z - \mu_i}{\sigma_i^2} + \frac{1}{B} \cdot \sum_{j=1}^{|\mathcal{S}|} \text{softmax}_j \cdot \frac{z - \mu_j}{\sigma_j^2}, \\
\frac{\partial \mathbb{D}_{KL}}{\partial \log\sigma_i^2} &= \frac{1}{B} \cdot \left( -\frac{1}{2} + \frac{(z - \mu_i)^2}{2\sigma_i^2} \right) + \frac{1}{B} \cdot \sum_{j=1}^{|\mathcal{S}|} \text{softmax}_j \cdot \frac{z - \mu_j}{\sigma_j^2} \cdot \epsilon \odot \sigma_i \odot 0.5.
\end{align}

$\square$

\subsection{Prior Parameter Non-Differentiability}

\textbf{Theorem 7} (Prior Parameter Non-Differentiability): The mixture prior parameters $\mu_j, \sigma_j^2, w_j$ from other clients do not receive gradients.

\textbf{Proof:} Prior parameters are fixed values from other clients' distributions, used as constants in the forward pass:

\begin{equation}
p_{\text{mixture}}(z) = \sum_{j=1}^{|\mathcal{S}|} w_j \cdot \mathcal{N}(z; \mu_j, \sigma_j^2 I),
\end{equation}

where $\mu_j, \sigma_j^2, w_j$ are fixed (not part of the computation graph for the current client).

In the gradient computation:

\begin{equation}
\frac{\partial \mathbb{D}_{KL}}{\partial \mu_j} = \frac{\partial}{\partial \mu_j} \mathbb{E}_{z \sim q_i} \left[ \log q_i(z) - \log p_{\text{mixture}}(z) \right] = 0,
\end{equation}

since $\mu_j$ is not a leaf node in the computation graph and autograd does not compute gradients for it.

$\square$
